\begin{boxedabstract}
Despite rapid advances in computer vision and generative AI, data sharing and the lack of standardized datasets remain major bottlenecks in medical imaging. Skin lesion datasets suffer from limited labeled segmentation data. We leverage HAM10000, a large public dermoscopy dataset, as our primary data source, to generate skin lesion images from binary masks. On this dataset, we train a ControlNet-based latent diffusion model that synthesizes photorealistic dermoscopic images conditioned on binary lesion segmentation masks. To enable fully automated generation, we additionally train a lightweight diffusion model that generates plausible lesion masks from scratch, forming an end-to-end pipeline capable of producing paired mask-image samples without any manual input. [Quantitative results to be added.]
\end{boxedabstract}

\begin{keywordsblock}
Dermoscopy, ControlNet, Diffusion Models, Data Augmentation, Skin Lesion Synthesis
\end{keywordsblock}

\vspace{2mm}
